\begin{figure}[!hbtp]
    \centering
    \resizebox{\columnwidth}{!}{%
    \begin{tikzpicture}[
        bloco/.style={
            rectangle,
            rounded corners=3pt,
            draw=black!70,
            fill=blue!10,
            line width=0.8pt,
            inner sep=6pt,
            align=center,
            font=\small,
            text width=2.8cm,
            minimum height=2.8cm
        },
        arrow/.style = {
            thick,
            ->,
            >={Stealth[length=2.5mm]}
        }
    ]

        \node (b1) [bloco] {\textbf{Bloco $N-1$}\\[3pt] \textit{Cabeçalho:}\\ \textit{hash} anterior\\ \textit{nonce}\\ raiz de Merkle\\[3pt] \textit{Transações}};
        \node (b2) [bloco, right=1.4cm of b1] {\textbf{Bloco $N$}\\[3pt] \textit{Cabeçalho:}\\ \textit{hash} anterior\\ \textit{nonce}\\ raiz de Merkle\\[3pt] \textit{Transações}};
        \node (b3) [bloco, right=1.4cm of b2] {\textbf{Bloco $N+1$}\\[3pt] \textit{Cabeçalho:}\\ \textit{hash} anterior\\ \textit{nonce}\\ raiz de Merkle\\[3pt] \textit{Transações}};

        \draw [arrow] (b2.west) -- node[above, font=\scriptsize\itshape] {\textit{hash}} (b1.east);
        \draw [arrow] (b3.west) -- node[above, font=\scriptsize\itshape] {\textit{hash}} (b2.east);

    \end{tikzpicture}%
    }
    \caption{Estrutura de uma cadeia de blocos: cada bloco contém um cabeçalho, com o \textit{hash} do bloco anterior e a raiz da árvore de Merkle, e o conjunto de transações. O encadeamento por \textit{hash} torna detectável a adulteração de blocos passados.}
    \label{fig:estrutura_blocos}
\end{figure}
